\section{Principios de diseño de gráficos}
\label{design_principles}
    Los principios que se aplicaron de forma explícita en cada figura provienen
    fundamentalmente de Few \cite{few2012show} y Tufte \cite{tufte2001visual}, los
    cuales son:

    \vspace{10pt}

    \begin{itemize}
        \item \textbf{Integridad gráfica:} la representación visual debe ser
        proporcional a las cantidades que representa, de modo que el factor de mentira
        se mantenga cercano a uno, por ello, las perspectivas tridimensionales, las
        sombras y el desplazamiento radial (\textit{explode}) deben evitarse, ya que
        distorsionan la magnitud visual de los datos
        \cite{few2012show}.

        \item \textbf{Razón dato-tinta:} la mayor parte del gráfico debe destinarse a
        representar datos evitando rejillas gruesas, marcos dobles, fondos y efectos
        decorativos que consumen la atención sin aportar información
        \cite{few2012show}.

        \item \textbf{Etiquetado completo:} todo gráfico debe ser autocontenido, con
        título informativo, ejes rotulados con unidades explícitas y cuando las
        categorías lo permiten, etiquetas numéricas directas \cite{few2012show}.

        \item \textbf{Uso funcional del color:} el color codifica información, no
        decora, puesto que los matices se reservan para las variables nominales, las
        escalas secuenciales para las cuantitativas y si solo hay una variable, se
        prefiere un color único \cite{few2012show}.

        \item \textbf{Escala honesta:} en los gráficos de barras el eje parte de cero,
        porque se compara por longitud; en los de dispersión es legítimo truncar el
        rango, porque se lee por posición relativa \cite{few2012show}.

        \item \textbf{Ordenamiento significativo:} cuando la categoría no tiene orden
        natural, organizarla por magnitud reduce el esfuerzo de comparación visual del
        lector \cite{few2012show}.

        \item \textbf{Jerarquía visual:} la cuadrícula y los ejes se dibujan por detrás
        de los datos y con menor peso visual, dejando la información en primer plano
        \cite{few2012show}.
    \end{itemize}